\documentclass[12pt]{article}
\usepackage[utf8]{inputenc}
\usepackage[russian]{babel}
\usepackage{amsmath}
\usepackage{amsfonts}
\usepackage{amssymb}

\title{Пример документа LaTeX}
\author{Автор}
\date{\today}

\begin{document}

\maketitle

\section{Введение}
Это введение к документу LaTeX. Здесь можно описать цель документа, контекст и основные положения, которые будут рассмотрены в тексте. LaTeX - это мощная система подготовки документов, особенно популярная в научной и технической литературе.

\section{Основная часть}
В основной части документа обычно размещается основной контент, анализ, формулы, таблицы и другие элементы, зависящие от темы документа. В этом примере основная часть остается пустой, так как основной акцент сделан на введение и заключение.

\section{Заключение}
В заключении подводятся итоги документа, делаются выводы и, при необходимости, формулируются рекомендации или направления для дальнейших исследований. Этот документ представляет собой базовую структуру для документа типа article с введением и заключением.

\end{document}